\subsection{Гистограмма выборки}

Введём следующее разбиение для данной выборки:
\[
    \tau=\{ 2.55, 2.70, 2.85, 3.00, 3.15, 3.30, 3.45 \}
\]

Обозначим полуинтервалы следующим образом:
\[
    \Delta_i=[\tau_i,\tau_{i+1}),\quad i=\overline{0,5}
\]

Для данного разбиения \( \tau \) имеем:
\[
    l_i=\abs{\Delta_i}=\tau_{i+1}-\tau_i=0.15,\quad i=\overline{0,5}
\]

Определим число элементов выборки \( \nu_i \), которые попали в \( \Delta_i \) интервал:

\begin{table}[h!]
\centering
\begin{tabular}{|c|c|c||c|c|c|}
\hline
$i$ & $\Delta_i$ & $\nu_i$ & $i$ & $\Delta_i$ & $\nu_i$ \\ \hline \hline
0 & $[2.55; 2.70)$ & 2 & 3 & $[3.00; 3.15)$ & 7 \\ \hline
1 & $[2.70; 2.85)$ & 0 & 4 & $[3.15; 3.30)$ & 3 \\ \hline
2 & $[2.85; 3.00)$ & 3 & 5 & $[3.30; 3.45)$ & 5 \\ \hline
\end{tabular}
\caption{Частотное распределение по интервалам $\Delta_i$}
\end{table}

Найдём высоту столбца \( h_i \) по формуле:
\[
h_i = \frac{\nu_i}{n\cdot l_i}
\]

Для наглядности приведём промежуточные значения в таблице:

\begin{table}[H]
\centering
\begin{tabular}{|c|c|c||c|c|c|}
\hline
$i$ & $\nu_i$ & $h_i$ & $i$ & $\nu_i$ & $h_i$ \\ \hline \hline
0 & 2 & 0.66667 & 3 & 7 & 2.33333 \\ \hline
1 &  0 & 0.00000 & 4 & 3 & 1.00000 \\ \hline
2 &  3 & 1.00000 & 5 & 5 & 1.66667 \\ \hline
\end{tabular}
\caption{Высоты столбцов гистограммы по интервалам \( \Delta_i \)}
\end{table}

Чтобы построить гистограмму, также воспользуемся библиотекой \textit{TiKZ} с~элементами автоматизации на \textit{Python}:

\begin{lstlisting}[caption=Функция построения гистограммы]
def generate_histogram():
  result = tikz_header
  for i in range(len(tau) - 1):
    result += f"\\addplot[fill=white, ybar interval, pattern={pattern_type[i % 2]}, draw=black, pattern color=black] coordinates {{"
    result += f"({tau[i]}, {heights[i]}) ({tau[i+1]}, {heights[i]})"
    result += "};\n"
  result += tikz_footer
  return result
\end{lstlisting}

Итого имеем:

\begin{figure}[h!]
\centering
\begin{tikzpicture}
\begin{axis}[
    width=0.9\textwidth, height=7cm,
    axis lines=middle,
    axis line style={-},
    xlabel={$x$},
    ylabel={$h$},
    xlabel style={right},
    ylabel style={above},
    ymin=0, ymax=3,
    xmin=2.5, xmax=3.5,
    ytick={0,0.5,...,3.5},
    xtick={2.55,2.70,2.85,3.00,3.15,3.30,3.45},
    xticklabels={2.55,2.70,2.85,3.00,3.15,3.30,3.45},
    grid=major,
    grid style={dashed, gray!50},
    ybar,
]
\addplot[fill=white, ybar interval, pattern=dots, draw=black, pattern color=black] coordinates {(2.55, 0.66667) (2.7, 0.66667)};
\addplot[fill=white, ybar interval, pattern=north east lines, draw=black, pattern color=black] coordinates {(2.7, 0.0) (2.85, 0.0)};
\addplot[fill=white, ybar interval, pattern=dots, draw=black, pattern color=black] coordinates {(2.85, 1.0) (3.0, 1.0)};
\addplot[fill=white, ybar interval, pattern=north east lines, draw=black, pattern color=black] coordinates {(3.0, 2.33333) (3.15, 2.33333)};
\addplot[fill=white, ybar interval, pattern=dots, draw=black, pattern color=black] coordinates {(3.15, 1.0) (3.3, 1.0)};
\addplot[fill=white, ybar interval, pattern=north east lines, draw=black, pattern color=black] coordinates {(3.3, 1.66667) (3.45, 1.66667)};
\end{axis}
\end{tikzpicture}
\caption{Гистограмма выборки}
\end{figure}

Проверка достоверности полученного результата:
\[
    S_{\text{гист.}}=\sum_{i=0}^5 h_i\cdot l_i=\dfrac{3}{20}\cdot\left(\dfrac{2}{3}+0+1+\dfrac{7}{3}+1+\dfrac{5}{3}\right)=\dfrac{3}{20}\cdot\dfrac{20}{3}=1\quad\square
\]